\scanpage{36}\mspage{4}

En d'autres termes, en termes du groupe $H$, et des deux sous-groupes $U_0$, $U_1$ de $H$ et de leurs structures de $\Lambda$-modules inversibles $L_0$, $L_1$ (automatique si $\Lambda$ est un anneau absolu $\mathbb Z$ ou $\mathbb Z/n\mathbb Z$), on trouve

a) L'accouplement $L_0\otimes L_1\simeq\Lambda$, par la condition que $u_0\otimes u_1=-1$ sss

\[
\begin{aligned}
 &u_1u_0u_1\in\operatorname{Transp.}(U_0,U_1).\quad\text{ou encore}\\
 &u_1u_0u_1\in\operatorname{Transp}(U_0,U_1)\cap\operatorname{Transp}(U_1,U_0)\overset{\mathrm{def}}{=}N^-
\end{aligned}
\]

\margin{ou encore $u_0\in U_0^*$ et \deleted{$(u_1u_0u_1)($} $\operatorname{int}(u_1u_0u_1)(u_0)=u_1$ i.e. (38) ci-dessous}

b) L'application $\sigma_0:U_0^*\longrightarrow N^-$ par la condition $\sigma_0(u_0)=u_1u_0u_1$, où $u_1\in U_1$ est l'unique élément de $U_1$ satisfaisant à la condition que $u_1u_0u_1\in N^-$, et alors on a

\[
\text{c)}\qquad u_1=\sigma_0(u_0)(u_0)=\underbrace{u_1u_0u_1}u_0\underbrace{u_1^{-1}u_0^{-1}u_1^{-1}}
\]

i.e.

\[
 u_0u_1u_0u_1^{-1}u_0^{-1}u_1^{-1}=1\quad\text{i.e }\ u_0u_1u_0=u_1u_0u_1\tag*{(38)}
\]

formule qui est équivalente encore à $u_0\otimes u_1=-1$. En termes de $e_0,e_1$, (38) s'écrit aussi

\[
\begin{aligned}
 e_0(\beta)e_1(\gamma)e_0(\beta)=e_1(\gamma)e_0(\beta)e_1(\gamma)&\qquad\text{sss }\beta\gamma=-1\\
 &(\beta\in L_0^*,\ \gamma\in L_1).
\end{aligned}\tag*{(39)}
\]

La relation (38) \deleted{est} entre $U_0$ et $U_1$ est le graphe d'une bijection entre $U_0^*=e_0(L_0^*)$ et $U_1^*=e_1(L_1^*)$, qu'on peut noter

\[
u_0\longmapsto u_0^\vee:U_0^*\longrightarrow U_1^*
\]

et satisfaisant

\[
(\deleted{\unsure{\lambda}}u_0^\lambda)^\vee=(u_0^\vee)^{\lambda^{-1}},\qquad u_0\in U_0^*,\ \lambda\in\Lambda^*\tag*{(40)}
\]

en notant exponentiellement l'opération \inserted{externe} de $\Lambda$ sur $U_0,U_1$.

% archive sheet number: 35
